\chapter{Introduction}
\label{ch:intro}

\section{Motivation}

Institutional data is rich, but only a few people can ask questions of it. A student wants to know
in which courses their attendance is short, a head of department wants the five best students in
their department, and an accounts officer wants the fee demands still unpaid. Each of these is one
SQL query away, yet in practice each depends on a screen someone built in advance or on a request to
the administrators. Large language models (LLMs) promise to remove this bottleneck: the user types a
question and the model writes the query. This is the text-to-SQL problem, and benchmarks such as
Spider \citep{yu2018spider} and BIRD \citep{li2023bird} show that models have become capable at it.

\section{The Problem: Capable Is Not Safe}

Accuracy on a benchmark is not what stops an ERP from leaking data. Benchmarks measure whether the
generated query returns the right answer and assume that the asker may see everything. A real
deployment needs the opposite guarantee: whatever the model writes, \emph{the asker must never
receive data outside their rights}. Three properties of LLMs make this hard.

\begin{itemize}
  \item \textbf{The output is untrusted.} A model can hallucinate tables, ignore an instruction, or
        be steered by a malicious question (prompt injection, \citealp{perez2022ignore}) into
        producing a query the operator never intended.
  \item \textbf{Permissions are about people, not queries.} A student may read their own payments but
        nobody else's; a head of department sees only their department. The same English question
        must give different, correctly scoped answers to different users.
  \item \textbf{Production databases are large.} A university ERP has hundreds of columns across many
        schemas, which a model cannot read in full and which a policy must still cover completely.
\end{itemize}

Prompting the model to ``only show the user their own data'' is a request, not a control. The
guarantee has to be enforced by code that does not depend on the model behaving.

\section{NatQuery in Brief}

\emph{NatQuery} is a Python library and command-line tool that turns English questions into
PostgreSQL queries and runs them only after they have been checked against the asker's rights. Its
single design rule is:

\begin{quote}
\emph{The language model proposes SQL. NatQuery decides whether, and in what form, it runs.}
\end{quote}

The model never connects to the database. Each proposed query is parsed into a syntax tree and
verified against a role-specific view of the schema, so tables and columns the role may not see do
not exist for it. Mandatory row filters are then attached, sensitive columns are masked, and only
the rewritten query is executed, on a read-only connection with time and row limits. A query that
breaks a rule is returned to the model with a precise reason so that it can repair it; a query that
cannot be repaired is refused. Who the user is comes from a pluggable authentication layer, what each
role may do is declared in one policy file, and every step is written to a tamper-evident,
hash-chained audit log. For databases spread over several schemas, NatQuery also passes the user's
identity to PostgreSQL row-level security, so that the database enforces the same scopes as a
second, independent line of defence.

\section{Contributions}

The contributions of this work are:

\begin{enumerate}
  \item A \textbf{security-first text-to-SQL architecture} in which LLM output is never executed
        directly: parse, role-view check, row filtering, masking, repair loop and audit
        (Chapters~\ref{ch:arch} and~\ref{ch:how}).
  \item A \textbf{policy-driven access model} with pluggable authentication, role inheritance, row
        filters bound to user attributes, multi-schema support, and database row-level security as a
        second enforcement layer (Chapters~\ref{ch:how} and~\ref{ch:features}).
  \item A \textbf{three-part evaluation}: execution accuracy on the public BIRD Mini-Dev benchmark,
        an adversarial security test on AdventureWorks, and a purpose-built synthetic university ERP
        (87 tables, 549 columns) with 40 role-specific questions at three difficulty levels
        (Chapter~\ref{ch:bench}).
  \item An \textbf{open test bed}: the ERP schema, seed data, policy and question set are released with
        the code, together with a no-LLM verification pass that separates access-layer errors from
        model errors.
\end{enumerate}

Table~\ref{tab:glance} summarises the results that the rest of the report supports.

\begin{table}[htbp]
\centering
\caption{NatQuery at a glance: the three evaluation tracks and the code base.}
\label{tab:glance}
\small
\renewcommand{\arraystretch}{1.2}
\begin{tabularx}{\textwidth}{L{2.6cm} Y Y}
\toprule
\hdr \textbf{Track} & \textbf{What is tested} & \textbf{Headline result} \\
\midrule
BIRD Mini-Dev & 500 questions, 11 databases, PostgreSQL, security layer on, open model \code{gpt-oss-120b} &
  45.6\% official execution accuracy (50.6\% with a numerically tolerant scorer); median 1.08\,s and \$0.0004 per question \\
AdventureWorks security & 3 roles, 227 hostile SQL queries (560 runs), 40 auth/audit checks, 30 natural-language attacks &
  0 leaks, 0 writes; 511 of 511 must-refuse queries refused; an unprotected text-to-SQL leaked on 7 of 8 comparable attacks \\
Synthetic university ERP & 10 schemas, 87 tables, 549 columns, 10 roles, row-level security; 40 questions in three levels &
  51 of 51 database checks; 41 of 41 no-LLM pipeline checks; 28 of 40 questions correct (70\%); all 5 access-control questions safe \\
Code base & about 5,300 lines of Python in 57 modules &
  855 tests, about 97\% branch coverage; 96 of 97 manual mutants killed \\
\bottomrule
\end{tabularx}
\end{table}

\section{Scope and Limitations}

NatQuery targets PostgreSQL and read-only questions. Its accuracy depends on the language model
behind it; the results reported here use an open-weight model on a free API tier, with limited
prompt size, and should be read as a lower bound. The ERP test bed is synthetic and its policy was
written by the same authors as the questions, so it shows that the system handles a database of this
size and shape, not that any particular real policy is correct. These limits are discussed in
Chapter~\ref{ch:bench} and Chapter~\ref{ch:conclusion}.

\section{Organisation of the Report}

Table~\ref{tab:organisation} maps the chapters to their content.

\begin{table}[htbp]
\centering
\caption{Organisation of the report.}
\label{tab:organisation}
\small
\renewcommand{\arraystretch}{1.2}
\begin{tabularx}{\textwidth}{L{2.2cm} Y}
\toprule
\hdr \textbf{Chapter} & \textbf{Content} \\
\midrule
2 & Architecture diagram, components and trust boundaries \\
3 & How NatQuery works: request path, enforcer, repair loop, schema selection, multi-schema and row-level security, audit \\
4 & Features: security, access control, accuracy, operations, extensibility and quality \\
5 & Benchmarking: BIRD Mini-Dev, AdventureWorks security, synthetic ERP \\
6 & Conclusion, use cases, limitations and future work \\
Appendices & Policy and schema reference, the 40 ERP questions with results, reproduction \\
\bottomrule
\end{tabularx}
\end{table}
